\section{The Kepler problem}

In this section we prove that the Kepler problem on the Heisenberg group is not integrable.
This result will follow as a corollary from a more general theorem.
Our proof is based on the Morales--Ramis theorem \cite{Morales:01::b1} which gives
the following necessary conditions for the integrability.
\begin{Theorem} \label{thm:MoRa}
 If a Hamiltonian system is Liouville integrable, with meromorphic
 first integrals, then the identity component of the differential
 Galois group of variational equations along any non-constant
 solution is Abelian.
\end{Theorem}

Two remarks are in order before we jump to application. One is the
``fine print'' of the above theorem: if the variational equation is
not Fuchsian, then only rational first integrals can be~treated. This
will turn out to be the case here, due to the choice of the particular
solution.

Secondly, because the Hamiltonian~\eqref{KH} itself is not
meromorphic, thanks to algebraic potential
\begin{equation*} %\label{potkep}
 V= - \frac{\kappa}{\sqrt{\big(x^2+y^2\big)^2+{16}z^2}},
\end{equation*}
a slight modification to Theorem~\ref{thm:MoRa} is necessary. We describe it in Appendix~\ref{App0}.

When applying the above, the main difficulty is connected
with determination of properties of the differential Galois group of
variational equations. If they can be reduced to a second order equation, then we can use the decisive Kovacic
algorithm~\cite{3}. Sometimes it is possible to~show that the
considered variational equations contain as a subsystem an equation
for which the differential Galois group is known, e.g., hypergeometric
equation and its confluent form. Here~we give a very useful example
which we will apply later.
\begin{Theorem}[H.P.~Rehm, 1979] \label{thm:par}
 Assume that complex parameters $\alpha\neq 0$, $\beta$, and $\gamma$
 of the parabolic cylinder equation
 \begin{equation} \label{eq:5}
 w''(z) - \big(\alpha^2 z^2 +2 \alpha \beta z + \gamma\big) w(z)=0
 \end{equation}
 are such that $\big(\beta^2-\gamma\big)/\alpha$ is not an odd integer. Then
 its differential Galois group is $\mathrm{SL}(2,\mathbb{C})$.
\end{Theorem}
This theorem was proved in \cite{Rehm} and later in~\cite{baby} it was
proved in another way with the help of the Kovacic algorithm.
